%-----------------------------------------------------------------------------%
\section{Conclusion}
\label{bab:5}
%-----------------------------------------------------------------------------%

This chapter summarizes the findings and proposes future research directions based on the evaluation results. Section 5.1 outlines the conclusion, reflecting on how the resource defragmentation and network-aware policies address the core research questions. This chapter concludes with Section 5.2, which proposes future works and open challenges to further enhance the descheduling architecture.

\subsection{Conclusion}
\label{sec:conclusion}

This research designed and evaluated a multi-level Kubernetes Descheduler architecture that combines request-based resource defragmentation, actual-usage pre-eviction protection, and network-latency-aware eviction filtering.

\textbf{Identifying fragmented Nodes.} \textit{ResourceDefragmentation} identifies fragmented Nodes directly from their request-space resource shape, without a utilization threshold. It selects drain candidates from average utilization and CPU-memory bin quality, orders them by the Resource Imbalance Index (RII) and Free-Space Index (FSI), and then evicts the Pod whose predicted landing Node, found through \textit{predictSchedulerTarget}, has the best projected bin score. A target ceiling and pack-upward guard prevent overfilling a Node or relocating a workload to an emptier one. This single predicted-target criterion is sufficient to drive effective packing and defragmentation.

\textbf{Effectiveness against the baseline.} ResourceDefragmentation reclaimed four Nodes in S1, three in S2, and one in S4, all cases where \textit{HighNodeUtilization} performed no eviction. It cut stranding by roughly \(74\%\) in S2 (\(3.727 \to 0.971\)) and \(66\%\) in S4 (\(3.243 \to 1.112\)), and by \(91\%\) in the heterogeneous-drain S5 (\(1.338 \to 0.123\)) while reclaiming the source Node with only two evictions against the baseline's three. It also preserved or increased balanced headroom in every fragmentation scenario, every run converged to a zero-eviction pass within three passes, and no run left a Pod Pending.

\textbf{Actual-usage filtering and application protection.} Our evaluation demonstrates that runtime usage can be incorporated into the Kubernetes Descheduler as a pre-eviction safety gate without modifying the candidate ordering of the active balancing strategy. The \textit{Actual Usage Evictor} blocked the busy API Pod under \textit{ResourceDefragmentation}, \textit{HighNodeUtilization}, and \textit{LowNodeUtilization}, while idle candidates remained eligible. Over a common \(30\,s\) post-event window, the control failure rates ranged from \(16.67\%\) to \(17.24\%\), and each control experienced approximately \(4.8\,s\) of service disruption. Every treatment recorded \(0\%\) failure and no disruption. The placement cost varied by strategy. ResourceDefragmentation retained an eviction but lost the \(0.816\) stranding reduction produced by moving the API Pod. HighNodeUtilization reclaimed two rather than three source Nodes. LowNodeUtilization substituted an idle fallback and retained the same three evictions and \(85.7\%\) memory-deviation reduction. These results show that runtime protection can preserve application availability while allowing continued Descheduler progress when a compatible idle candidate remains available.

\textbf{Network-aware filtering and latency trade-offs.} Our evaluation demonstrates that network-aware filtering can be effectively integrated into the Kubernetes Descheduler via a dual-layer architecture. We developed a system combining a global pre-eviction filter (\textit{NetworkCostEvictor}) with strategy-level assessments, using real-time telemetry or topology labels to intercept latency-degrading evictions. Our analysis indicates that this architecture preserves communication efficiency without compromising resource-balancing objectives, provided non-network-sensitive pods are available to sacrifice. Specifically, our results show that the system capped latency at $\sim$1.2ms (load-balancing) and $\sim$6.1ms (multi-dimensional skews) while preventing all cross-region pod spills, avoiding the $\sim$82ms to $\sim$98ms latency penalties observed in upstream baselines. However, our evaluation of the \textit{HighNodeUtilization} scenario also revealed a critical systemic limitation: isolated descheduler intent can be subverted by the network-blind default Kubernetes scheduler, which frequently misrouted evicted pods to degraded nodes based on raw compute availability. This confirms that high-precision network locality requires explicit synchronization with the primary cluster scheduler. Ultimately, our findings prove that network protection and geometric resource balance can be achieved together when alternative eviction candidates exist; otherwise, administrators must strictly prioritize between full cluster defragmentation and network locality preservation. 

Overall, the results support a simple conclusion: selecting eviction candidates by the balance of their predicted landing Nodes is enough to drive useful Node reclamation and dimensional defragmentation when combined with request feasibility, bounded utilization, pack-upward behavior, actual-usage filtering, and network-aware safety checks.

\section{Suggestions}
\label{sec:future_works}

\begin{itemize}
    \item \textbf{Broader Resource Dimensions:} Extend the framework to additional resource dimensions, including storage and GPU, so that target prediction and bin-quality scoring can handle workloads beyond CPU and memory.
    \item \textbf{Scheduler Coordination:} Establish and evaluate scheduler coordination mechanisms that expose the predicted target, or a compatible scoring signal, to the kube-scheduler, ensuring that predicted-target placements align with actual scheduling decisions and eliminating the misalignment that currently bounds maximum defragmentation precision.
    \item \textbf{Pending-Pod Awareness:} Introduce Pending-pod awareness into the eviction pipeline so that drain candidates are selected with knowledge of already-unschedulable workloads, ensuring evictions free the right resource type on the right Node to unblock Pending Pods rather than delaying them further.
    \item \textbf{Threshold Sensitivity Analysis:} Conduct a sensitivity analysis of the drain threshold, target ceiling, and related parameters across various workloads and cluster types to characterize their effect on node reclamation, stranding reduction, and convergence stability.
    \item \textbf{Broader Generalizability Evaluation:} Evaluate the plugin across broader cluster shapes, node heterogeneity, and workload mixes to assess generalizability beyond the five scenarios studied.
    \item \textbf{Traffic-Volume Weighting:} Weight network costs using service-mesh traffic volumes so that heavily used dependency paths receive stronger protection.
\end{itemize}
